\slajdid{06-s002}
\begin{frame}[label=06-s002]
\gusttekst\setlength{\parskip}{3pt}
Pored informacija o vrednostima prikazanih brojno, digitalni sistem obrađuje i druge informacije.

Te informacije su takođe predstavljen kombinacijama logičkih nula i jedinica odnosno „kodovane“ su.

Na primer svaki procesor sadrži neku vrstu registra\\u kojem se čuvaju informacije o izvršenim aritmetičkim operacijama

Program status register (PSR) - Flag register - Condition code register (CCR)

Na primer informacija je kodovana u bitima
\begin{quote}\upshape
C – Carry flag – rezultat sabiranja, oduzimanja; može da se tretira i kao borrow\par
Z – Zero flag – rezultat je nula\par
S/N – Sign, Negative flag – rezultat je negativan\par
V/O/W – Overflow flag – rezultat operacije u drugom komplenetu ima overflow\par
\ldots
\end{quote}
Sam „hardver“ ne vodi računa ako je nastupilo prekoračenje opsega.\\
Ali na višem „programskom“ nivou korisnik može da proverava informaciju, kodovanu u PSR, da li je prilikom aritmetičkih operacija nastupilo prekoračenje i da, ako smatra da su potrebne, preduzme odgovarajuće akcije.
\end{frame}
